After the meeting, I began to contemplate the product design of this device. My starting point was to design a very non-intrusive device, meaning this device would lean towards being a companion with simple and limited interaction. Given that smartphones and personal computers are already effective communication tools in our daily lives, I did not intend to add complex or high-density communication functionalities to this device, as it could not fully replace our phones or computers in terms of functionality and usability. Therefore, I aimed to introduce differentiated features that would make it unique and enjoyable to use, aligning with the theme of "enchanted everyday objects." With these preliminary ideas, I participated in another meeting to elaborate on our device's feature set. This meeting focused primarily on the theme of "enchantment." Considering that a snow globe might not be an everyday object for everyone, we revised our choice to a lamp and developed a feature set, building our personas and storybook around that.

The feature set includes:
\begin{enumerate}
    \item The device should support multiple users and connect them.
    \item The device should be able to display the ambient lighting conditions of other devices.
    \item Each user should be represented by a small semi-transparent figure, with the ambient lighting showing as the figure lighting up or not.
    \item The figure should be interactive, such as tapping it to send a signal.
    \item The device should also display light signals representing the weather conditions of other users, including temperature and weather.
    \item The device should support sending and receiving voice messages.
\end{enumerate}

We build our scenarios, storyboards and customer journey maps based on these discussions.

These features were discussed as a group with the preference of all members. If I were fully in charge of the design process, I would propose some changes:

\begin{enumerate}
    \item The device should not support more than two users, as this would significantly limit the interaction design by shifting from one-to-one to multi-to-multi interactions. Additionally, multi-to-multi interactions could potentially introduce social tension among users, which I would personally prefer to avoid.
    \item The voice message function should be removed. In my view, if we cannot achieve a complete and highly usable voice message function, we should instead focus on simpler but more fun features.
\end{enumerate}

Over the weekend, I organized a small focus group session with some former colleagues from my time working at the design center of WeChat, including two software engineers and one designer. I introduced them to our design and sought their opinions. To avoid bias, I refrained from arguing or commenting on their views, providing only a neutral description of the product. They provided some feedback which prompted further thought. The result was used during the low fidelity prototype phrase.

Additionally, at this current point in time when I organise this document, one reflection I have is that our design didn't focus deeply enough on understanding how to achieve the "enchanted" aspect. My initial thought was something "wonderful" or "magical", but as we delved into the literature, we found that "enchanted" is more often defined as \cite{mccarthy2006experience}:
"\textit{For an individual consumer, enchantment is an experience as deeply felt yet fleeting set of emotional commitments involving wonderment, anticipation of joy, euphoria, and an expanded sense of human potential}” 
Enchantment plays a big role in the design and adoption of connectivity devices, it engages closely and intimately with the object of affective attachment, absorbing its characteristics . 
That said, by being enchanted, the object may not be as rational as I thought \cite{belk2021enchantment}. I could have design the product in a more  "absurd" and "incredible" way and focus more on emotions. 